\chapter{Changing experiments at runtime}\label{ch:runtime}

Most properties of an experiment can be changed after it is loaded, for
example to study the sensitivity as a function of the baseline or the running
time. After any change, call \code{simulate!} again so that the simulated data
match the new experiment.

\section{Exposure}\label{sec:exposure}

\begin{apifun}{target\_mass(e)\qquad target\_mass!(e, mass)}
The fiducial detector mass (in the units of the AEDL file, usually kt).
\end{apifun}
\begin{apifun}{running\_time(e; flux=1)\qquad running\_time!(e, years; flux=:all)}
The running time of one flux, or setting it for one or all fluxes.
\end{apifun}
\begin{apifun}{source\_power(e; flux=1)\qquad source\_power!(e, power; flux=:all)}
The source power: beam power, or useful parent decays per year.
\end{apifun}

\noindent Since the event numbers are proportional to all three, a scan over
the exposure is
\begin{julia}
for years in (1, 2, 5, 10)
    running_time!(nova, years / 2)            # half neutrinos, half antineutrinos
    simulate!(truth)
    println(years, "  ", chi2(OscParams(truth; δ=0.5π)))
end
\end{julia}

\section{Baseline and matter profile}

\code{baseline!(e, L)} changes the baseline and recomputes the density profile
as the AEDL file defines it, \code{density\_profile!(e, lengths, densities)}
sets any profile (Section~\ref{sec:profiles}). Note that the flux decreases as
$1/L^2$ when the baseline changes:
\begin{julia}
nf = Experiment("NFstandard.glb")
for L in 1000:1000:8000
    baseline!(nf, L)
    simulate!(truth)
    println(L, "  ", total_events(nf; rule=1))
end
\end{julia}

\section{Systematics and energy windows}

\code{systematics!}, \code{chi2\_function!} and \code{systematic\_errors!}
(Chapter~\ref{ch:chi2}) change the systematics treatment of rules, and
\code{energy\_window!} their energy window.

\section{AEDL variables}

AEDL files can use variables that are defined by the program. They are passed
as keyword arguments when the experiment is loaded; numbers become variables
and vectors become lists:
\begin{julia}
nf = Experiment("NFvar.glb"; BASELINE=3000, emax=50)
\end{julia}
In the AEDL file, they are used like any variable,
\code{\$baseline = BASELINE} (Section~\ref{sec:aedl-variables}). Lists are
named with a leading \code{\%} in AEDL; the keyword is given without it.
This is the way to change things that cannot be changed after loading, such
as the number of bins or the energy range.

\section{The low-pass filter\texorpdfstring{$^\ast$}{*}}

The filter of Section~\ref{sec:filter} is set in AEDL with
\code{\$filter\_state} and \code{\$filter\_value}; at runtime, use
\code{glbSetFilterStateInExperiment(i, GLB\_ON)} and
\code{glbSetFilterInExperiment(i, σ)} with the experiment number
\code{i} (Chapter~\ref{ch:compat}).
